% Member 4 - final report, IV.F Comparison with a Published Feature-Space Baseline
\subsection{Comparison with a Published Feature-Space Baseline}
\label{sec:mixstyle}

MixStyle~\\cite{zhou2021mixstyle} is the closest published competitor of this
family: it removes the style of the input, and it does so by mixing the
channel-wise statistics of the feature maps with those of another sample.
It therefore addresses the same problem as style augmentation, but in
feature space and without generating any image. We include it in
Table~\ref{tab:main_results} as the reference point against which the
generative approach has to justify itself.

\paragraph{How it is run.}
MixStyle changes the architecture, so it is a separate group with its own
configuration fingerprint, and it is trained on the \emph{original} Photo
images: the layer does its work inside the network, so there is no
augmentation set to build and no manifest to check. Its hyper-parameters
are the published defaults ($p = 0.5$, $\alpha = 0.1$), and the layer is
switched off at evaluation time, as the method prescribes. The optimiser,
the step budget, the seeds and the evaluation splits are those of every
other group, so the comparison is a comparison of the training data and
the architecture, not of the training protocol.

\paragraph{The result.}
MixStyle reaches a cross-domain average of $42.2\pm2.6\%$, against
$40.8\pm2.3\%$ for the Original group: a paired difference of $+1.37$
points ($t = -1.84$, $p = 0.140$), which is not significant. The styled
groups of Table~\ref{tab:main_results} reach $47.2$ to $47.4\%$ over the
same five seeds. The generative augmentation is therefore ahead of the
feature-space baseline by about five points, and that gap is large compared
with the seed-to-seed spread of either method.

\paragraph{What the comparison does and does not show.}
It shows that in this setting the generative approach is the more effective
of the two. It does not show that MixStyle is a weak method in general. Two
limitations of the comparison should be stated. First, we ran MixStyle with
its published defaults and did not tune it for PACS; a tuned feature-space
method might do better, and we report this baseline as a reference point,
not as a refutation. Second, the two methods are not interchangeable: the
feature-space layer adds parameters to the network, whereas the generative
method leaves the classifier untouched and produces inspectable images.
The comparison we can make is therefore about the accuracy achieved under
one shared protocol, and that is how we present it.

\paragraph{The instructive part.}
The more interesting observation is what MixStyle does \emph{not} change.
The decomposition of Section~\ref{sec:main_results} attributes the gain of
the styled groups to the amount of style they apply and to nothing else.
MixStyle removes style in a completely different way --- in feature space,
and without changing the amount of appearance variation in the training
images --- and it recovers only a fraction of the gain. This is consistent
with the picture that emerges from the rest of the study: the training
images themselves have to carry the appearance variation, and it is the
variation that has to be \emph{visible} to the classifier, not merely
removed from the representation.
